\subsection{0+1维动力学:算符方法}
在 1.1.2 节中我们已经建立了态的幺正演化 $|\psi(t)\rangle = U(t,0)|\psi(0)\rangle$，并由此导出了 Schr\"{o}dinger 方程。现在我们系统地比较两种等价的动力学描述方式。
\topictitle{Schr\"odinger 绘景}
在 Schr\"{o}dinger 绘景中，\textbf{态随时间演化，算符不变}。具体地：
\begin{itemize}
\item 态矢量携带全部时间依赖性：
\begin{equation}
|\psi(t)\rangle_S = e^{-iHt}\,|\psi(0)\rangle
\end{equation}
\item 可观测量算符 $\mathcal{O}_S$ 不含时间（除非有显式的时间依赖）。
\item 期望值为
\begin{equation}
\langle \mathcal{O} \rangle(t)
= {}_S\langle\psi(t)|\,\mathcal{O}_S\,|\psi(t)\rangle_S
\end{equation}
\end{itemize}
这是最直观的描述：物理系统的``状态''随时间改变，而``仪器''（算符）保持不动。Schr\"{o}dinger 方程
\begin{equation}
i\frac{d}{dt}|\psi(t)\rangle_S = H\,|\psi(t)\rangle_S
\end{equation}
就是这一绘景中的运动方程。
\topictitle{Heisenberg 绘景}
在 Heisenberg 绘景中，\textbf{态不变，算符随时间演化}。两种绘景通过幺正变换联系：
\begin{equation}
|\psi\rangle_H \equiv |\psi(0)\rangle_S\,,\qquad
\mathcal{O}_H(t) \equiv e^{iHt}\,\mathcal{O}_S\,e^{-iHt}
\end{equation}
由于 $U(t)=e^{-iHt}$ 是幺正的，期望值在两种绘景中完全相同：
\begin{equation}
{}_S\langle\psi(t)|\,\mathcal{O}_S\,|\psi(t)\rangle_S
= \langle\psi(0)|\,e^{iHt}\,\mathcal{O}_S\,e^{-iHt}\,|\psi(0)\rangle
= {}_H\langle\psi|\,\mathcal{O}_H(t)\,|\psi\rangle_H
\end{equation}
这保证了两种绘景的物理等价性——\color{blue}所有可观测的概率和期望值完全一致。\normalcolor Heisenberg 绘景中算符的运动方程可以通过对 $\mathcal{O}_H(t)$ 求时间导数得到：
\begin{align}
\frac{d}{dt}\mathcal{O}_H(t)
&= \frac{d}{dt}\bigl(e^{iHt}\,\mathcal{O}_S\,e^{-iHt}\bigr) \notag\\
&= iH\,e^{iHt}\,\mathcal{O}_S\,e^{-iHt}
- e^{iHt}\,\mathcal{O}_S\,e^{-iHt}\,iH \notag\\
&= i[H,\,\mathcal{O}_H(t)]
\end{align}
即
\begin{equation}
\frac{d}{dt}\mathcal{O}_H(t) = i[H,\,\mathcal{O}_H(t)]
\end{equation}
这是 \textbf{Heisenberg 运动方程}。

\begin{center}\vspace{2ex}
\textbf{***}
\vspace{2ex}\end{center}

两种绘景在数学上严格等价，选择哪一种纯粹是方便性的问题：

\begin{itemize}
\item Schr\"{o}dinger 绘景更接近非相对论量子力学的直觉——波函数 $\psi(\pmb{x},t)=\langle\pmb{x}|\psi(t)\rangle_S$ 随时间演化。

\item Heisenberg 绘景在相对论量子场论中具有根本性的优势。在这种绘景中，场算符 $\phi_H(\pmb{x},t)$ 同时依赖于空间和时间，因此可以被视为时空点 $x^\mu=(\pmb{x},t)$ 的函数——这使得 Lorentz 协变性得以显式体现\footnote{这也是为什么在第 4 章讨论散射理论时，``入''态和``出''态被定义为不随时间改变的 Heisenberg 绘景态。}。
\end{itemize}
下面讨论\textbf{相互作用绘景}（Dirac 绘景）是这两种绘景的混合体;但此绘景本身存在些许问题。
\topictitle{相互作用绘景和Dyson级数}
在Schr\"odinger和Heisenberg绘景中，系统的动力学演化都可以被精确处理；而在现实情况中，动力学行为十分复杂，基本没有解析解，所以在这种情况下我们将完整的哈密顿量$H$分为自由部分$H_0$相互作用部分$V$：
\begin{equation}
    H=H_0+V
\end{equation}
同时约定由自由部分的$H_0$承担时间演化，
\begin{equation}
    \begin{aligned}
        \text{态矢量的演化}:&\ket{\psi(t)}_{ip}=e^{iH_0t}\ket{\psi(t)}_S=e^{iH_0t}e^{-iHt}\ket{\psi}\\
        \text{算符的演化}:&\mathscr{O}_{ip}(t)=e^{iH_0t}\mathscr{O}e^{-iH_0t}
    \end{aligned}
\end{equation}其中约定：$\mathscr{O}=\mathscr{O}(0)_H=\mathscr{O}_S$$\ket{\psi}=\ket{\psi}_H=\ket{\psi(0)}_S$
\vspace{2ex}
    \hrule
\vspace{2ex}\par
幺正演化和模长(期望)是等价的，所以：
\begin{equation}
    \begin{aligned}
        {}_{ip}\bra{\psi(t)}\mathscr{O}_{ip}(t)\ket{\psi(t)}_{ip}=&\bra{\psi}e^{iHt}\underbrace{e^{-iH_0t}e^{iH_0t}}_\mathbf{1}\mathscr{O}\underbrace{e^{-iH_0t}e^{iH_0t}}_\mathbf{1}  e^{-iHt}\ket{\psi}\\
        &=
        \bra{\psi}e^{iHt}\mathscr{O}  e^{-iHt}\ket{\psi}\\
        &=
        \bra{\psi(t)}\mathscr{O} \ket{\psi(t)}=\bra{\psi}\mathscr{O}_H(t) \ket{\psi}
    \end{aligned}
\end{equation}



态的演化是$\ket{\psi(t)}_{ip}=e^{iH_0t}e^{-iHt}\ket{\psi}$,可以推导时间演化算符，
\begin{equation*}
    \begin{aligned}
    \text{$t$时刻}:&\ket{\psi(t)}_{ip}=e^{iH_0t}e^{-iHt}\ket{\psi}\\
    \text{$t_0$时刻}:&\ket{\psi(t_0)}_{ip}=e^{iH_0t_0}e^{-iHt_0}\ket{\psi}\to \ket{\psi}=e^{iHt_0}e^{-iH_0t_0}\ket{\psi(t_0)}_{ip}
\end{aligned}
\end{equation*}
便可得到从$\ket{\psi(t)}_{ip}$演化到$\ket{\psi(t_0)}_{ip}$的方程，便定义时间演化算符：
\begin{equation}
    \begin{aligned}
        \ket{\psi(t)}_{ip}=&
        e^{iH_0t}e^{-iHt}\ket{\psi}\\
        &=
        e^{iH_0t}e^{-iHt}e^{iHt_0}e^{-iH_0t_0}\ket{\psi(t_0)}_{ip}\\&=
        \underbrace{e^{iH_0t}e^{-iH(t-t_0)}e^{-iH_0t_0}}_{U(t,t_0)}\ket{\psi(t_0)}_{ip}
    \end{aligned}\qquad\text{且}\quad
    U(t,t)=1
\end{equation}
现在来推导时间演化算符$U(t,t_0)$的形式，对$U(t,t_0)=e^{iH_0t}e^{-iH(t-t_0)}e^{-iH_0t_0}$两端求时间的导数：
\begin{equation*}
        \begin{aligned}
        \frac{d}{dt}U(t,t_0)=&\frac{d}{dt}\left\{e^{iH_0t}e^{-iH(t-t_0)}e^{iH_0t_0}\right\}\\
        =&iH_0U(t,t_0)-ie^{iH_0t}He^{-iH(t-t_0)}e^{iH_0t_0}\\
        =&iH_0U(t,t_0)-ie^{iH_0t}H\underbrace{e^{-iH_0t}e^{-iH_0t}}_{\mathbf{1}}e^{-iH(t-t_0)}e^{iH_0t_0}\\
        =&iH_0U(t,t_0)-ie^{iH_0t}He^{-iH_0t}U(t,t_0)\\
        =&iH_0U(t,t_0)-ie^{iH_0t}(\underbrace{H_0}_\text{自由哈密顿量时间演化后不变}+V)e^{-iH_0t}U(t,t_0)\\
        =&-ie^{iH_0t}Ve^{-iH_0t}U(t,t_0)\\
        =&-iV(t)U(t,t_0)
    \end{aligned}
\end{equation*}
对上式积分:
\begin{equation*}
    \begin{aligned}
        &\int^{t}_{t_0}dU(t,t_0)=\int^{t}_{t_0}dt\left[-iV(t)U(t,t_0)\right]\\
        &\rightarrow U(t,t_0)-U(t_0,t_0)=-i\int^{t}_{t_0}dt\left[V(t)U(t,t_0)\right]\\
        &\rightarrow U(t,t_0)=1-i\int^{t}_{t_0}dt\left[V(t)U(t,t_0)\right]\\
    \end{aligned}
\end{equation*}
 不断迭代这个积分,可以得到对于$V$的幂次展开的积分式：
\[    \begin{aligned}
        \rightarrow U(t,t_0)=1-&i\int^{t}_{t_0}dt_1V(t_1)+(-i)^2\int^{t}_{t_0}dt_1\int_{t_0}^{t_1}dt_2V(t_1)V(t_1)\\
        &+(-i)^3\int^{t}_{t_0}dt_1\int^{t_1}_{t_0}dt_2\int^{t_2}_{t_0}dt_3V(t_1)V(t_2)V(t_3)+\cdots
    \end{aligned}\]
取渐近极限$t\rightarrow\infty$和$t_0\rightarrow-\infty$,就给出了时间演化算符的渐近展开形式：
\begin{equation}\label{dyson}
    \begin{aligned}
        \rightarrow U(\infty,-\infty)=1-&i\int^{\infty}_{-\infty}dt_1V(t_1)+(-i)^2\int^{\infty}_{-\infty}dt_1\int_{-\infty}^{t_1}dt_2V(t_1)V(t_1)\\
        &+(-i)^3\int^{\infty}_{-\infty}dt_1\int^{t_1}_{-\infty}dt_2\int^{t_2}_{-\infty}dt_3V(t_1)V(t_2)V(t_3)+\cdots
    \end{aligned}
\end{equation}
这里利用算符的\textbf{编时乘积}(按时间排序)写为编时乘积的形式，编时算符的形式为：
\[
\begin{aligned}
T\{V(t)\} &= V(t), \\
T\{V(t_1)V(t_2)\} &= \theta(t_1-t_2)V(t_1)V(t_2)
+ \theta(t_2-t_1)V(t_2)V(t_1)\\
&\ldots
\end{aligned}
\]
其中 $\theta(\tau)$ 是阶跃函数，$\tau>0$ 时等于 $+1$ 而 $\tau<0$ 时等于 $0$。
$n$ 个 $V$ 的编时乘积是对全部 $n!$ 个 $V$ 的置换的求和，
每一个都对所有的 $t_1,\ldots,t_n$ 积分，所以方程(\refeq{dyson})可以写为
\begin{equation*}
    \begin{aligned}
        U(\infty,-\infty)
        = 1 + \sum_{n=1}^{\infty} \frac{(-i)^n}{n!}
        \int_{-\infty}^{\infty} dt_1\,dt_2 \cdots dt_n\,
        T\{V(t_1)\cdots V(t_n)\}
\end{aligned}
\end{equation*}
如果不同时刻的所有 $V(t)$ 是对易的，
这个级数可以进行求和，其和是
\[
\begin{aligned}
U = \exp\!\left(-i\int_{-\infty}^{\infty} dt\,V(t)\right).
\end{aligned}
\]
当然，通常不是这种情况；一般而言，上式甚至不收敛，
它充其量是 $V$ 中出现的任何耦合常数的渐近展开\footnote{关于渐近展开的入门介绍，Erdelyi的短篇专著《Asymptotic Expansions》（Dover, 1956）非常有用}式。
然而在一般情况下，上式有时写为
\begin{equation}\label{dyson_ser}
        U = T \exp\!\left(-i\int_{-\infty}^{\infty} dt\,V(t)\right)   
\end{equation}

\subsection{0+1 维动力学: 泛函方法}

前文中算符方法处理了量子动力学；本节引入一种概念上截然不同但数学上严格等价的描述——\textbf{路径积分}
与算符方法不同，路径积分的出发点不是 Hilbert 空间中的算符和态矢量，而是经典构型空间中所有可能路径的加权求和——
每一条路径 $q(t)$ 被赋予权重 $e^{iS[q]}$，其中 $S[q]$ 是沿该路径计算的经典作用量

\medskip
量子力学的全部运动学结构，归根到底建立在两类对易关系之上。
\begin{subequations}
  \begin{align}
    &[Q_i,Q_j]=0\\
    &[P_i,P_j]=0\\
    &[Q_i,P_j]=\delta_{ij}
  \end{align}
\end{subequations}
其中对易的算符提供了本征关系,以及完备正交基
\begin{subequations}\label{Eigen_relation_1}
    \begin{align}
      &[Q_i,Q_j]=0\to Q\ket{q}=q\ket{q}\quad\text{以及}
      \left\{\begin{aligned}
        &\braket{q|q'}=\delta(q-q')\\
        &\mathbf{1}=\int dq \ket{q}\bra{q}
      \end{aligned}\right.\\
      &[P_i,P_j]=0\to P\ket{p}=p\ket{p}\quad\text{以及}
      \left\{\begin{aligned}
        &\braket{p|p'}=\delta(p-p')\\
        &\mathbf{1}=\int dp \ket{p}\bra{p}
      \end{aligned}\right.
    \end{align}
\end{subequations}
而不对易的算符提供了：
\begin{equation}\label{Eigen_relation_2}
  [Q_i,P_j]=\delta_{ij}\to \braket{q|p}=\frac{1}{2\pi}e^{ipq}
\end{equation}
\medskip
路径积分的自然语言是 Heisenberg 绘景——在这种描述中，\textbf{所有的时间演化都被编码在算符 $e^{iHt}$ 中}。
Heisenberg 绘景中的位置和动量算符定义为：
\begin{subequations}\label{eq:Heis-ops}
    \begin{align}
        &Q(t) = e^{iHt}\,Q\,e^{-iHt}\\
        &P(t)= e^{iHt}\,P\,e^{-iHt}
    \end{align}
\end{subequations}
而本征态是：
\begin{subequations}
  \begin{align}
    &\ket{q,t}=e^{iHt}\ket{q}\\
    &\ket{p,t}=e^{iHt}\ket{p}
  \end{align}
\end{subequations}
现在来稍加证明：对于Heisenberg绘景下的算符和态，依然满足本征关系(\refeq{Eigen_relation_1})以及(\refeq{Eigen_relation_2})
\[\begin{aligned}
  Q(t)\ket{q,t}&
  =\left(e^{iHt}\,Q\,e^{-iHt}\right)\left(e^{iHt}\ket{q}\right)\\
  &=e^{iHt}\,Q\,\ket{q}=q\,e^{iHt}\ket{q}=q\ket{q,t}
\end{aligned}\]
当然，对于算符$P(t)$和以及$\ket{p,t}$是同样的；所以本征关系、完备正交基以及(\refeq{Eigen_relation_2})依然是正确的：
\begin{subequations}
  \begin{align}
    &\braket{q,t|q',t}=\delta(q-q')\quad \mathbf{1}=\int dq \ket{q,t}\bra{q,t}\\
    &\braket{p,t|p',t}=\delta(p-p')\quad \mathbf{1}=\int dp \ket{p,t}\bra{p,t}\\
    &\braket{q,t|p,t}=\frac{1}{2\pi}e^{ipq}
  \end{align}
\end{subequations}
\topictitle{传播子：核心动力学}
现在考虑两个\textbf{不同时刻}的 Heisenberg 位置本征态的内积：
\begin{equation}\label{eq:propagator}
  K(q_f, t_f;\, q_i, t_i)
  \;\equiv\; \langle q_f, t_f\,|\,q_i, t_i\rangle
  \;=\; \langle q_f|\,e^{-i(t_f - t_i)H(Q(t),P(t))}\,|q_i\rangle
\end{equation}
传播子 (\ref{eq:propagator}) 包含算符 $e^{-i(t_f-t_i)H(Q(t),P(t))}$，对于是算符的哈密顿量无法直接计算。
处理的办法取一个小间隔的传播子，令$t_f+\epsilon=t$、$t_i=t$，其中$\epsilon$是一个无穷小量，而这一间隔的传播子是：
\[
\braket{q+dq,t+\epsilon|q,t}=\braket{q+dq|e^{-i\epsilon H(Q(t),P(t))}|q}
\]
利用Taylor展开：
\[\braket{q+dq,t+\epsilon|q,t}=\braket{q+dq|1-i\epsilon H(Q(t),P(t))+\mathscr{O}(\epsilon^2)+\dots|q}\]
\color{blue}这里利用Weyl排序，让哈密顿量的$Q(t)$都在$P(t)$的左边，对于非Weyl序的，可以利用对易关系替换\normalcolor；在展开式中插入完备基$\mathbf{1}=\int dp\ket{p}\bra{p}$\,(规定p的积分区间是$(-\infty,\infty)$):
\[\braket{q+dq,t+\epsilon|q,t}=\int dp \,
\bra{q+dq}\left\{1-i\epsilon H(Q(t),P(t))+\mathscr{O}(\epsilon^2)+\dots\right\}\ket{p}\braket{p|q}
\]
这种情况下,可以利用本征关系，哈密顿量从一个算符变为一个“数”
\[\braket{q+dq,t+\epsilon|q,t}=\int dp \,
\bra{q+dq}\left\{1-i\epsilon H(q,p)+\mathscr{O}(\epsilon^2)+\dots\right\}\ket{p}\braket{p|q}
\]
由于$\epsilon$是无穷小量，自然的上式能会到指数形式：
\begin{equation}\label{q+dq}
  \begin{aligned}
    \braket{q+dq,t+\epsilon|q,t}=&\int dp \,
  \bra{q+dq}\left\{1-i\epsilon H(q,p)+\mathscr{O}(\epsilon^2)+\dots\right\}\ket{p}\braket{p|q}\\
    &=\int dp \,e^{-i\epsilon H(q,p)}\braket{q+dq|p}\braket{p|q}\\
    &=\int dp \,e^{-i\epsilon H(q,p)}\frac{1}{2\pi}e^{ip(q+dp)}e^{ipq}\\
    &=\int \frac{dp}{2\pi} \,e^{-i\epsilon H(q,p)+ip\,dq}\\
    &=\int \frac{dp}{2\pi} \,e^{-i \epsilon\left(H(q,p)+p\,\frac{dq}{\epsilon}\right)}
  \end{aligned}
\end{equation}
回到任意间隔，将时间区间 $[t_i, t_f]$ 等分为 $N$ 份，每份宽度为
$$\varepsilon = \frac{t_f - t_i}{N+1}$$
对于从初态到末态的传播子，对每一个时间间隔$\epsilon$的态$\ket{q_k,t_k}$都插入完备基：
\begin{equation}
  \langle q_f, t_f\,|\,q_i, t_i\rangle=\int dq_1 \dots dq_N\braket{q_f,t_f|q_N,t_N}\Braket{q_N,t_N|q_{N-1},t_{N-1}}\dots \braket{q_1,t_1|q_i,t_i}
\end{equation}
其中每一个间隔的传播子都可以利用(\refeq{q+dq})代替，所以：
\begin{equation}
   \braket{q_f,t_f|q_i,t_i}=\int \left(\prod_{k=1}^{N}dq_k\right)\left(\prod_{k=0}^{N}\frac{dp_k}{2\pi}\right)\exp[i\epsilon\sum_{k=1}^{N}\left\{H(q_k,p_{k-1})+p_k\,\frac{dq_k}{\epsilon} \right\}]
\end{equation}
安全的取极限$\epsilon\to 0$,并约定记号$\epsilon=d\tau$所以
\begin{equation}
   \braket{q_f,t_f|q_i,t_i}=\lim_{\epsilon \to 0}\int \left(\prod_{k=1}^{N}dq(\tau)\right)\left(\prod_{k=0}^{N}\frac{dp(\tau)}{2\pi}\right)\exp[i\int \left\{H(q(\tau),p(\tau))+p_k(\tau)\,\dot{q(\tau)} d\tau\right\}]
\end{equation}
上式称为相空间中的路径积分，它对坐标空间中所有满足边界条件 $q(t_f) = q_f$ 和 $q(t_i) = q_i$ 的路径 $q(t)$ 进行积分，也对动量空间中所有的 $p(t)$ 路径进行积分
 
\topictitle{经典极限与驻相近似}
路径积分提供了理解量子-经典过渡的最直观方式。暂时恢复 $\hbar$：
\[
  K = \int\mathcal{D}q\;\exp\!\left\{\frac{i}{\hbar}S[q]\right\}
\]
当 $\hbar \to 0$ 时，相位 $S/\hbar$ 对路径的微小偏离变化极为剧烈，
不同路径的贡献因快速振荡而\textbf{相消干涉}——
恰如集团分解论证所利用的 Riemann-Lebesgue 机制。
唯一的例外是\textbf{驻相点}（使相位变化率为零的路径）：
\begin{equation}
  \frac{\delta S}{\delta q(t)}\bigg|_{q = q_{\text{cl}}} = 0
  \qquad\Longleftrightarrow\qquad
  \frac{d}{dt}\frac{\partial L}{\partial\dot{q}}
  - \frac{\partial L}{\partial q} = 0
\end{equation}
这正是经典力学的 \textbf{Euler--Lagrange 方程}。因此：
经典力学中粒子沿唯一的使作用量取极值的经典轨迹运动；
量子力学中粒子"探索"所有可能的路径，
但在 $\hbar\to 0$ 极限下只有经典路径附近的路径彼此相干增强。

将 $q(t) = q_{\text{cl}}(t) + \eta(t)$（$\eta(t_i) = \eta(t_f) = 0$）代入作用量，
展开到涨落 $\eta$ 的二阶\footnote{
一阶项 $\delta S/\delta q|_{q_{\text{cl}}} = 0$ 因经典方程成立而消失。}：
\begin{equation}
  S[q] = S[q_{\text{cl}}]
  + \frac{1}{2}\int\!dt\,dt'\;\eta(t)\,
    \frac{\delta^2 S}{\delta q(t)\delta q(t')}\bigg|_{q_{\text{cl}}}\,\eta(t')
  + O(\eta^3)
\end{equation}
对涨落 $\eta$ 做高斯积分，得到半经典（WKB）近似：
\begin{equation}
  K \;\approx\; \mathcal{N}\;\exp\!\left\{\frac{i}{\hbar}S[q_{\text{cl}}]\right\}
\end{equation}
其中预指数因子 $\mathcal{N}$ 由二阶涨落算符的函数行列式决定。

